% answers_short.tex —— 2006 年·普通高等学校招生考试·上海卷理科 · 简短答案速查
% 由 tools/make_short_answers.py 自 answers.tex 的 \gksolhead/\gkanswer 提取，
% 并与 2006/2006_数学真题_上海理_答案要点.md 逐题比对后生成，勿手改（改源 MD 后重跑）。

\begin{gknotice}
\textbf{使用说明}：本节只给最终答案便于快速核对，完整推导见解析版。答案为编写者独立推导，\textbf{非官方标准答案}。
\end{gknotice}

\gkpart{填空题}
\begin{itemize}[leftmargin=2.6em]
  \item[\textbf{1.}] $m\in\{-\sqrt3,1,\sqrt3\}$
  \item[\textbf{2.}] $\dfrac1{\sqrt2}$
  \item[\textbf{3.}] $a=\dfrac12$
  \item[\textbf{4.}] $\dfrac16$
  \item[\textbf{5.}] $z=-1+i$
  \item[\textbf{6.}] $\dfrac{2\sqrt6}{5}$
  \item[\textbf{7.}] $\dfrac{x^2}{16}+\dfrac{y^2}{4}=1$
  \item[\textbf{8.}] $5$
  \item[\textbf{9.}] $\dfrac1{35}$
  \item[\textbf{10.}] $36$
  \item[\textbf{11.}] $k=0,\ |b|<1$
  \item[\textbf{12.}] $a\le10$
\end{itemize}

\gkpart{选择题}
\begin{center}\begin{tabular}{ccccc}
\toprule
题号 & 13 & 14 & 15 & 16 \\
\midrule
答案 & C & A & A & D \\
\bottomrule
\end{tabular}\end{center}

\gkpart{解答题}
\begin{itemize}[leftmargin=2.6em]
  \item[\textbf{17.}] 值域 $[-2,2]$，最小正周期 $\pi$
  \item[\textbf{18.}] 北偏东 $71^\circ$
  \item[\textbf{19.}] (1) $2$；(2) $\arccos\dfrac1{2\sqrt2}$
  \item[\textbf{20.}] (1) 原命题为真；(2) 逆命题为假命题
  \item[\textbf{21.}] (1) 公比为 $a$；(2) $b_n=1+\dfrac{n-1}{2k-1}$；(3) $k=2,3,4,5,6,7$
  \item[\textbf{22.}] (1) $b=9$；(2) 递减区间 $(-\infty,-c^{1/4}]$、$(0,c^{1/4}]$；递增区间 $[-c^{1/4},0)$、$[c^{1/4},+\infty)$；(3) 可推广为 $x^m+a/x^m$；单调区间见解析。$F$ 的最大值为 $(9/2)^n+(9/4)^n$，最小值为 $2^{n+1}$
\end{itemize}
